\setupinteraction[title={\getvariable{meta}{title}}]

\setuplayout[
   backspace=1cm, % Space on the left margin
   width=fit, % Adjust width to fit remaining space
   height=fit,
   topspace=0.5cm, % Top margin
   header=0.5cm, % Header size
   rightedge=0.5cm,
   leftedge=1cm,
   footer=1cm, % Footer size
   bottomspace=1cm, % Bottom margin
   rightmargin=3cm, % Space for margin notes
   rightmargindistance=0.5cm,
]

\setupbodyfont[mainface,9pt]

\setuphead[chapter]
          [
			header=empty,
            style=\itd,
            number=no,
            % before={\blank[big,force]},
            after={\blank[2*big,force]},
          ]

% \showgrid
\showlayout
\showframe

\starttext

\startbodymatter
\chapter[title={Lernzettel -
Stochastic},reference={lernzettel---stochastic}]

\section[title={Wahrscheinlichkeiten},reference={wahrscheinlichkeiten}]

\startnarrower[left]\setupindenting[-\leftskip,yes]

\margintext{{\em Zuffalsexperiments:} Ziehen, mit Zurücklegen, von zwei Kugeln aus
einer Urne mit fünf roten und drei blauen Kugeln.

{\em Ergebnismenge:} $S=\{rr; rb; br; bb\}$.

Das {\em Ereignis} $E = \{rr; rb; br\}$ bedeutet, dass mindestens eine
Kugel rot ist. Das {\em Gegenereignis} von $E$ ist $\overline{E}$
}

Die {\bf Wahrscheinlichkeiten} aller Ergebnisse eines Zuffalsexperiments
sind Zahlen im interval $\[0; 1\]$ mit Summe $1$. Sie bilden die {\em
Wahrscheinlichkeitsverteilung}. Die Ergebnisse fässt man in der {\bf
Ergebnismenge} zusammen, eine {\em Teilmenge} davon ist ein {\bf
Ergebnis}.

\margintext{Als {\bf fair} bezeichnet man ein Spiel, bei dem der Erwartungswert für
den Gewinn null ist. Gewinn = Auszahlung - Einsatz
}

\stopnarrower

\startframedtipp
{\bf Definition:} Wenn jedem Ergebnis eines Zufallsexperiments ein
Zahlenwert zugeordnet wird, spricht man von einer {\bf Zufallsgröße}.
Die {\bf Wahrscheinlichkeitsverteilung} ener Zufallsrgöße $X$ ist eine
Tabelle, bei der jedem Wert $k$ von $X$ die Wahrscheinlichkeit $P(X=k)$
zugeordnet ist. Für eine Zufallsgröße $X$ mit den Werten
$x_1,x_2, ..., x_n$ heißt $\mu = x_1 \cdot
P(X=x_1) + x_2 \cdot P(X = x_2)... + x_n \cdot P(X = x_n)$ {\bf
Erwartungswert} von $X$. Er gibt an, welchen Mittelwert man bei
ausreichend großer Versuchsanzahl auf lange Sicht erwartet.

\stopframedtipp

\startsectionlevel[title={Zufallsgrößen},reference={zufallsgrößen}]

\startframedtext[width=\dimexpr\textwidth+\rightmarginwidth\relax,frame=off,offset=none,toffset=1mm] % This spans textwidth + margin width
% \startframedtext[width=\dimexpr\textwidth+\rightmarginwidth+\rightmargindistance\relax,frame=off] % This spans textwidth + margin width

\startplacetable[location=here,title={Größen}]
\setupTABLE[each][align={middle,lohi}]
\setupTABLE[c][5][align={flushleft}]
\setupTABLE[row][1][background=color,backgroundcolor=gray,align={flushleft,lohi}]
\setupTABLE[row][2][background=color,backgroundcolor=gray]
\bTABLE[option=stretch,offset=1mm]
\bTR
	\bTD[nc=2] Basisgröße \eTD
	\bTD[nc=2] Basiseinheit \eTD
\eTR
\startCSV
Name, Symbol, Formel, Beschreibung
Varianz, $V$, $(x_1-\mu)^2\cdot P(X=x_1)+... +(x_n-\mu)^2\cdot P(X=x_n)$, Ein Maß für die Streuung
Standardabweichung, $\sigma$, $\sqrt{V}$, Abweichung der Werte von dem eigenen Durchschnittswert
Erwartungswert, $\mu$, $x_1 \cdot P(X = x_1) + x_2 \cdot P(X = x_2) + ... + x_n \cdot P(X = x_n)$
\stopCSV

\eTABLE
\stopplacetable

\stopframedtext% }}}

\section[title={Bedingte
Wahrscheinlichkeiten},reference={bedingte-wahrscheinlichkeiten}]

{\bf Beispiel am Urnenmodell}

In einer Urne sind {\bf 10} Kugeln, {\bf 5} davon sind Markiert
(Ereignis $M$). Also $P(M)=\frac{5}{10}=50\%$. Es gibt allerding {\bf
drei} von {\bf vier} roten Kugeln, welche Markiert sind und und {\bf
zwei} von {\bf sechs} nicht rote Kugeln. Wenn man nun beim ziehen vorher
schon weiß, welche Farbe die Kugel hat, bevor man die Markierung sieht,
verändert sich die Wahrscheinlichkeit auf $\frac{3}{4}=75\%$.

\margintext{Das bedingende Ereignis $R$ wird als Index notiert. Man liest $P_R(M)$:
\quotation{Wahrscheinlichkeit von $M$ unter der Bedingung $R$}
}

Man bezeichnet die Wahrscheinlichkeit für eine Markierung ($M$) unter
der Bedingung rot ($R$) als {\bf bedingte Wahrscheinlichkeit} und
schreibt.

\startformula
P_R(M) = \frac{P(M \cap R)}{P(R)}
\stopformula


\startframedtipp
Laplace-Versuch: Ein versuch, bei dem jedes Ergebnis\sidenote{Ergebnis ist jede Möglichkeit beim einmaligen ausführen eines Zufallsexperiments. Zum Beispiel bei einem Würfel ist jede Seite in Ergebnis. Ein Ereignis setzt sich aus ein oder mehrern Ergebnissen zusammen.} die gleich Wahrscheinlich ist
\stopframedtipp

\stopbodymatter

% \startbackmatter
% \component c_definitions
% \component c_bibliography
% \stopbackmatter

\stoptext
